\section{Discussion and Limitations}
\label{sec:discussion}

The framework fits the post-Grenfell regulatory direction directly. Safety cases
under the Building Safety Act ask duty-holders to demonstrate that risks are
understood and managed on an ongoing basis \citep{grenfell2024phase2,bsa2022};
a posterior risk distribution with explicit credible intervals and
missing-evidence flags is a more faithful object for that purpose than a
categorical score, and the expected-value-of-inspection machinery turns the same
posterior into a prioritisation rule for where to spend finite inspection effort.
The results support the premise: the occurrence model is calibrated on held-out
years and recovers the deprivation gradient; the consequence models are calibrated
out-of-time and isolate a clean detection$\to$spread signal; and the simulator
reproduces the observed spread orderings while exposing an archetype-specific
mechanism (detection helps flats and houses, but not mobility-limited care homes)
that a purely statistical model would miss; and the value-of-targeting study puts a
conservative lower bound of a roughly ninefold return on that prioritisation over
untargeted allocation. Taken together they move FRA from a static compliance
checklist \citep{pas79_1_2020,pas79_2_2020} toward an updateable, evidence-based
risk picture --- with the fully sequential, inspection-assimilating version
specified in Section~\ref{sec:model} as the programme's next stage.

The real-building case study sharpens what the archetypes are for. Grounding the
high-rise-flat archetype against Grenfell Tower shows that its scale parameters ---
storeys and occupant load --- matter far less to modelled risk than its regime
assumptions: doubling the storeys shifts the exposed-household failure probability by
only 2\%, whereas losing compartmentation and abandoning stay-put shifts expected
fatalities by two orders of magnitude. The archetype is a sound geometric proxy
precisely because it already encodes the safety-critical single-stair topology; its
unavoidable simplification is that, like every archetype, it presumes the design
strategy holds. The same study marks the framework's outer boundary: an interior
room--corridor--stair graph cannot represent external cladding-driven spread, the
mechanism that drove Grenfell, and must defer to the external-wall and safety-case
regime of the Building Safety Act for that hazard. A real-geometry v2 pipeline ---
replacing archetype graphs with per-building layouts drawn from the downloadable
building-specification sources we catalogue (\texttt{docs/building\_specs\_sources.md})
--- is the natural next step, but the external-wall hazard will remain a regulated,
per-building duty rather than something a population-scale model should attempt.

Several limitations bound these claims, and we have kept them visible rather than
smoothing them over. This is an England-only residential MVP built on archetype
graphs, not real building geometry; the simulator is a transparent stochastic
caricature, not fire-physics CFD, and its multi-storey houses over-predict spread
beyond the floor of origin because of an idealised open internal stair; the
care-home archetype's building-level evacuation-failure rate overstates the
observed communal-living casualty rate by roughly $1.4\times$ per occupant ---
and far more ($\sim$$4.8\times$) at the building-level any-failure measure ---
reflecting its no-staff-assist, stay-put-free full-evacuation assumption. Both
headline posteriors are exact full-data objects --- the consequence models by
NUTS on all 451{,}172 records, the occurrence model by PSIS-corrected Laplace
importance sampling on all 202{,}530 training rows ($\hat{k}=0.33$;
Section~\ref{sec:model}) after full-data NUTS failed under three
parameterisations --- each corroborated by subsample NUTS, ADVI and MLE
cross-checks; the occurrence model's linear time trend under-predicts the
steeper 2024/25 decline. The consequence models are
covariate-adjusted associations, not causal effects: the casualty--alarm odds ratio
falls below one purely through occupancy confounding, which is precisely why the
model carries a latent-state layer, but that layer de-confounds by assumption and
structure rather than by identification, and this paper reports the adjusted
associations that feed it rather than a validated causal estimate. The FRS
response-time covariate carries little trustworthy independent signal: it is
ecological (identified at FRS$\times$year resolution and entangled with the secular
trend), and its effect on spread is not even robust in sign --- below one on the
subsample, above one at full data --- so we retain it only as a control and report no
directional response-time effect. The exact full-data fit exposes a further honest
gap: it undercovers per-service held-out counts (an FRS-overdispersion a single
random intercept misses), so per-service predictive intervals should be read as
anti-conservative and a random-slope extension is future work. Data linkage is imperfect --- record-level
fires resolve only to FRS area, not LSOA; covariates mix vintages (IMD 2025, Census
2021, EPC 2012--2026); and EPC attributes enter only as LSOA aggregates, consistent
with the address-level redistribution restriction of Section~\ref{sec:data}. Finally,
fire is a rare event, so building-level validation is not possible from open data: our
calibration is at the aggregate (LSOA-year and FRS-area) level, and the worked
building-level decomposition of Section~\ref{sec:results-synthesis} is illustrative of
the composition, not a validated per-building prediction.

\paragraph{Equity of targeted prevention.} Because fire risk is strongly
deprivation-linked, any risk-targeted allocation concentrates visits in deprived
areas. We regard this as the intended behaviour of a prevention programme ---
resources flow to where expected harm is greatest --- but it carries obligations
the framework should make explicit rather than bury in a score: targeting
criteria must be auditable (a virtue of interpretable rate ratios over black-box
scores), the uncertainty attached to any area's ranking should be visible to
decision-makers, and allocations should be monitored for unintended effects, such
as areas of genuinely high but poorly evidenced risk (wide posteriors) being
systematically deprioritised relative to areas that are merely well measured.
The expected-benefit policy of Section~\ref{sec:results-decision} partially
addresses the latter by valuing uncertainty reduction, but a full fairness
analysis of risk-targeted fire prevention --- who gains, who is overlooked, and
under which valuation of harm --- is future work we consider necessary before
operational deployment.

\paragraph{The physics programme.} The simulator's mechanism-informed status is a
staging point, not an endpoint. A companion work builds the corresponding physics
engine (NIST CFAST zone models over the same archetypes, with era-derived boundary
materials and design fires keyed to the national ignition-source distribution) and
calibrates its parameters against the observed English spread margins with a
model-discrepancy term. Its headline is a caution this paper's framing anticipates:
the calibrated physics reproduces the aggregate spread surface, but aggregate
margins identify only one physical parameter (effective door-openness) --- the
remainder require \emph{timed}, scenario-resolved calibration targets, which is a
further argument for the record-level data access discussed above.

\paragraph{Secondary findings with policy implications.} Two supporting analyses
sharpen the operational picture. The no-alarm spread penalty is not uniform across
households --- it is roughly twice the pooled value for lone pensioners
(Section~\ref{sec:results-interactions}) --- giving the targeting layer a concrete,
data-backed reason to up-weight areas rich in that group beyond what their fire rate
alone implies, precisely where consequence decouples from occurrence, while the night
penalty is a whole-stock effect that warrants no dwelling-specific multiplier. And the
ignition-source analysis surfaces a data-infrastructure gap: England's central
incident taxonomy cannot isolate lithium-ion e-bike and e-scooter fires, the hazard
the sector is most anxious about, so the most actionable recommendation there is not a
model change but a taxonomy one (Section~\ref{sec:results-ignition}).

\paragraph{The inspection layer, meeting data.} The sharpest test of the paper's
central claim --- that an FRA is a noisy observation, not a risk score --- is the first
instantiation of the inspection layer on real assessments
(Section~\ref{sec:results-latent}): the apparatus runs end to end on 1{,}441 published
assessments and a single inspection demonstrably narrows a block's fire-rate posterior,
yet on one cross-sectional wave the latent's coupling to fire rate is a
remediation-confounded null, so a snapshot FRA is \emph{not} a validated proxy for
persistent fire propensity. That is exactly the situation the framework is designed for
--- inspections assimilated as noisy evidence under explicit uncertainty --- and it
names the missing ingredient: a second assessment wave to turn the retrospective check
into a prospective test and point-identify the assessor-noise parameters.
